% !TeX root = ../drafts/a-ring-switching.tex
\section{Ring switching}\label{annex:rs}

Turning $64$ evaluation claims over $\E$ into one opening of a commitment whose values are $64$ bits wide. This annex is the general construction; \S\ref{sec:ringswitch} instantiates it on the packed Flock witness $\qflock$.

\subsection{Setting}\label{rs:setting}

The two fields are those of \S\ref{sec:vm}: $\K=\F_2[x]/(x^{64}+x^4+x^3+x+1)$ inside $\E=\K[y]/(y^3+y+1)$. As $\F_2$-vector spaces, $\K$ has basis $x^0,\dots,x^{63}$, and $\E$ has the $192$ products
\[
e_w=x^iy^j\in\E,\qquad w=64j+i,\quad 0\le i<64,\quad 0\le j<3.
\]

The witness is $64$ bit-valued multilinears $Q_0,\dots,Q_{63}:\cube{\nflock}\to\{0,1\}$. We commit to them packed, one $\K$ element per point:
\[
\qflock:\cube{\nflock}\to\K,\qquad \qflock(u)=\sum_{i=0}^{63}Q_i(u)\,x^i.
\]
The commitment scheme proves statements of one shape. Given a target $T\in\E$ and an MLE-friendly weight $W:\cube{\nflock}\to\E$ (Definition~\ref{def:ipcs}), an opening proves
\[
\sum_{u\in\cube{\nflock}}W(u)\,\qflock(u)=T.
\]
The weight is $\E$-valued, the committed values are $\K$-valued, and the sum is taken in $\E$.

What arrives is not of that shape. It is $64$ claimed values $s_i\in\E$, one per bit polynomial, at a point $r\in\E^{\nflock}$ they all share. An honest prover sends
\[
s_i=\mle{Q_i}(r)=\sum_{u\in\cube{\nflock}}\eq(r,u)\,Q_i(u),\qquad 0\le i<64,
\]
where $\eq$ is the equality indicator of \S\ref{sec:prelim}, here in its general form
\[
\eq:\E^{\nflock}\times\E^{\nflock}\to\E,
\qquad
\eq(a,b)=\prod_{n=1}^{\nflock}\bigl(a_nb_n+(1+a_n)(1+b_n)\bigr),
\]
and $\mle{Q_i}$ is the multilinear extension of $Q_i$, as the tilde marks every extension below.
Ring switching~\cite{DP26} turns those $64$ claims into a single statement the opening can prove.

\subsection{The reduction}\label{rs:reduction}

Fix an $\F_2$-linear map $\Phi:\E\to\E$:
\begin{equation}
\begin{split}
\sum_{u\in\cube{\nflock}}\Phi\bigl(\eq(r,u)\bigr)\,\qflock(u)
&=\sum_{i=0}^{63}x^i\sum_{u\in\cube{\nflock}}Q_i(u)\,\Phi\bigl(\eq(r,u)\bigr)\\
&=\sum_{i=0}^{63}x^i\,\Phi\Bigl(\sum_{u\in\cube{\nflock}}Q_i(u)\,\eq(r,u)\Bigr)
=\sum_{i=0}^{63}x^i\,\Phi(s_i).
\end{split}
\label{eq:complete}
\end{equation}
The reduction is:
\begin{enumerate}
\item once the $s_i$ are fixed, the verifier draws $\Phi$ at random;
\item it computes the target $T=\sum_{i=0}^{63}x^i\Phi(s_i)$;
\item the commitment scheme opens $\qflock$ against the weight $W:\cube{\nflock}\to\E$, $W(u)=\Phi\bigl(\eq(r,u)\bigr)$, with target $T$; if that opening accepts, the verifier accepts the $64$ claims.
\end{enumerate}
Everything now turns on which $\Phi$ we can afford to draw.

\subsection{Which maps are sound}\label{rs:sound-maps}

\paragraph{The error, transposed.} Suppose the $s_i$ are wrong. The errors $\delta_i=s_i-\mle{Q_i}(r)\in\E$, for $0\le i<64$, are not all zero, and they are fixed before $\Phi$ is drawn. By \eqref{eq:complete} the target the verifier computes differs from the honest one by $\sum_{i=0}^{63}x^i\Phi(\delta_i)$, so the opening succeeds only if
\[
\sum_{i=0}^{63}x^i\,\Phi(\delta_i)=0.
\]
Expand each error in the $\F_2$-basis of $\E$, say $\delta_i=\sum_{w=0}^{191}\delta_{i,w}e_w$ with $\delta_{i,w}\in\F_2$, and exchange the two sums:
\begin{equation}
\sum_{i=0}^{63}x^i\,\Phi(\delta_i)
=\sum_{i=0}^{63}x^i\sum_{w=0}^{191}\delta_{i,w}\,\Phi(e_w)
=\sum_{w=0}^{191}\Phi(e_w)\,\Delta_w,
\qquad
\Delta_w:=\sum_{i=0}^{63}\delta_{i,w}x^i\in\K.
\label{eq:transpose}
\end{equation}
The $x^i$ are an $\F_2$-basis of $\K$, so some $\Delta_w\neq0$. The check passes only if $\sum_{w=0}^{191}\Phi(e_w)\,\Delta_w=0$. No fixed $\Phi$ rules that out, so the draw must make it unlikely.

\paragraph{Linear maps are Frobenius sums.} Squaring is additive in characteristic $2$, so
\[
\Phi(a)=\sum_{k=0}^{191}c_k\,a^{2^k},\qquad a\in\E,\quad c_k\in\E,
\]
is $\F_2$-linear for any choice of coefficients. Every $\F_2$-linear map $\E\to\E$ arises this way, and exactly once. Distinct tuples give distinct maps: their difference would be a nonzero polynomial of degree at most $2^{191}$ with $2^{192}$ roots. There are $|\E|^{192}$ tuples and as many maps. Call $\supp=\{k:c_k\neq0\}\subseteq\{0,\dots,191\}$ the \emph{support} of $\Phi$. Substituting into the check above,
\begin{equation}
\sum_{w=0}^{191}\Phi(e_w)\,\Delta_w=\sum_{k\in\supp}c_k\,\pair_k,
\qquad
\pair_k:=\sum_{w=0}^{191}e_w^{2^k}\,\Delta_w\in\E.
\label{eq:pairing}
\end{equation}
The $\pair_k$ depend only on the error, the $c_k$ only on the randomness. So we need the map $\K^{192}\to\E^{|\supp|}$, $\Delta\mapsto(\pair_k)_{k\in\supp}$, to be injective: a nonzero error has to leave some $\pair_k\neq0$ for the coefficients to act on.

\paragraph{Sixty-four terms.} That map is $\K$-linear, since the $\Delta_w$ enter linearly and the Frobenius touches only the constants $e_w$. It sends $\K^{192}$, of $\K$-dimension $192$, into $\E^{|\supp|}$, of $\K$-dimension $3|\supp|$, so injectivity needs $|\supp|\ge64$. Take $\supp=\{0,\dots,63\}$ and suppose $\pair_k=0$ for all $k<64$. Two steps force $\Delta=0$.

First, split $\pair_k$ over $\K$. Write $w=64j+i$ and $y_k:=y^{2^k}$, so that $e_w^{2^k}=x^{i2^k}y_k^{\,j}$ and
\[
\pair_k=\sum_{j=0}^{2}y_k^{\,j}\,\pair_{j,k},
\qquad
\pair_{j,k}:=\sum_{i=0}^{63}\Delta_{64j+i}\,x^{i2^k}\in\K,\qquad 0\le j<3.
\]
Now $y_k\notin\K$. Squaring is a bijection of $\K$ onto itself, so $y^{2^k}\in\K$ would force $y\in\K$, whence $\E=\K[y]=\K$. Since $[\E:\K]=3$ is prime, any element outside $\K$ generates $\E$ over $\K$, so $\{1,y_k,y_k^2\}$ is a $\K$-basis of $\E$, and $\pair_k=0$ forces $\pair_{0,k}=\pair_{1,k}=\pair_{2,k}=0$.

Second, read the $64$ equations $\pair_{j,k}=0$, at fixed $j$, as a single polynomial. Put
\[
D_j(v)=\sum_{i=0}^{63}\Delta_{64j+i}\,v^i\in\K[v],
\qquad\text{so that}\qquad
\pair_{j,k}=D_j\bigl(x^{2^k}\bigr).
\]
Then $D_j$ has degree at most $63$ and vanishes at the $64$ points $x,x^2,x^4,\dots,x^{2^{63}}$. Those are distinct, because $x$ has degree $64$ over $\F_2$ and so has $64$ distinct conjugates. A nonzero polynomial of degree at most $63$ cannot have $64$ roots, so $D_j=0$ and every $\Delta_{64j+i}$ vanishes. A nonzero error therefore leaves some $\pair_k\neq0$.

The count is tight on both sides: $64$ terms give $64$ equations $\pair_k=0$ over $\E$, which is $192$ equations over $\K$, against exactly $192$ unknowns $\Delta_w\in\K$.

\subsection{The map we draw}\label{rs:map}

We need a random $\Phi$ of support $\{0,\dots,63\}$, cheap to apply. Sixty-four independent coefficients would need $64$ challenges and $64$ multiplications per application, and a false claim would then survive with probability $2^{-192}$. Six challenges are enough, at less than $2^{-160}$.

Draw $f_0,\dots,f_5\leftarrow\E$ once the $s_i$ are fixed, and compose six two-term maps: for an input $a\in\E$, put $a_0=a$ and
\begin{equation}
a_{p+1}=a_p+f_p\,a_p^{2^{2^{5-p}}}\in\E\quad(0\le p<6),
\qquad
\Phi(a)=a_6.
\label{eq:phi}
\end{equation}
Each stage is $\F_2$-linear, so $\Phi$ is. Write $k=\sum_{p=0}^{5}k_p\,2^{5-p}$ for the binary digits of $k$. The composition expands to
\begin{equation}
\Phi(a)=\sum_{k=0}^{63}c_k\,a^{2^k},
\qquad
c_k=\prod_{p\,:\,k_p=1}f_p^{\,2^{(k\bmod 2^{5-p})}}.
\label{eq:ck}
\end{equation}
Stage $p$ fires when $k_p=1$ and contributes a factor $f_p$, which every later stage that fires raises to its own Frobenius power; those shifts sum to $k\bmod 2^{5-p}$. Each $k$ from $0$ to $63$ occurs exactly once, so the $c_k$ are $64$ distinct monomials in $f_0,\dots,f_5$, one for every Frobenius power the previous section requires.

\paragraph{Soundness.} Fix a nonzero error before the challenges are drawn. The previous section gives some $\pair_k\neq0$, and distinct monomials cannot cancel, so
\[
\sum_{k=0}^{63}c_k\,\pair_k
\]
is a nonzero polynomial in $f_0,\dots,f_5$. Its total degree is at most the largest degree of a $c_k$, reached at $k=63$ where every bit is set:
\[
2^{31}+2^{15}+2^7+2^3+2+1<2^{32},
\]
so by Schwartz--Zippel (Lemma~\ref{lem:sz}) the check passes with probability less than $2^{32}/2^{192}=2^{-160}$, before the opening's own soundness error.

\subsection{Evaluating the weight}\label{rs:weight}

An opening reduces $\sum_{u\in\cube{\nflock}}W(u)\,\qflock(u)=T$ to a claim about $\qflock$ at a random point $r'\in\E^{\nflock}$, and the verifier must evaluate the multilinear extension of $W$ at $r'$, written $\mle{W}(r')$, by itself. Two facts turn that into a short formula.

\emph{Fact 1 ($\eq$ in characteristic $2$).} In each factor the two $a_nb_n$ terms cancel, identically in $a$ and $b$:
\[
\eq(a,b)=\prod_{n=1}^{\nflock}(1+a_n+b_n).
\]

\emph{Fact 2 (Frobenius passes through $\eq$ on the boolean hypercube).} Squaring respects sums and products and fixes bits, so for $u\in\cube{\nflock}$, by Fact~1,
\[
\eq(r,u)^{2^k}=\prod_{n=1}^{\nflock}(1+r_n+u_n)^{2^k}=\prod_{n=1}^{\nflock}\bigl(1+r_n^{2^k}+u_n\bigr)=\eq\bigl(r^{2^k},u\bigr),
\qquad r^{2^k}:=\bigl(r_1^{2^k},\dots,r_{\nflock}^{2^k}\bigr).
\]

Expand $W$ at $r'$ with \eqref{eq:ck} and Fact~2:
\[
\begin{aligned}
\mle{W}(r')&=\sum_{u\in\cube{\nflock}}W(u)\,\eq(r',u)
=\sum_{u\in\cube{\nflock}}\Phi\bigl(\eq(r,u)\bigr)\,\eq(r',u)\\
&=\sum_{k=0}^{63}c_k\sum_{u\in\cube{\nflock}}\eq\bigl(r^{2^k},u\bigr)\,\eq(r',u).
\end{aligned}
\]
The inner sum is $\eq(r^{2^k},r')$, since $\eq(u,\cdot\,)$ is its own multilinear extension for every $u$. By Fact~1 again,
\[
\boxed{\ \
\mle{W}(r') \;=\; \sum_{k=0}^{63}\, c_k \prod_{n=1}^{\nflock}
\bigl( 1 + r_n^{2^k} + r'_n \bigr).
\ \ }
\]

\subsection{Cost}\label{rs:cost}

One application of $\Phi$ takes $63$ squarings and $6$ multiplications: the Frobenius ladder climbs $32+16+8+4+2+1$ steps, and each stage multiplies once by its challenge. The same recurrence run on the coefficient table gives the $c_k$ of \eqref{eq:ck}: stage $p$ raises its $2^p$ entries to the $2^{2^{5-p}}$ power and multiplies each by $f_p$; the raising is $32$ squarings whatever $p$ is.

\begin{center}
\begin{tabular}{lrr}
& squarings & multiplications\\[2pt]
\hline
one application of $\Phi$ & $63$ & $6$\\
the target $T$, $64$ applications & $4032$ & $384$\\
the coefficients $c_k$, once & $192$ & $63$\\
the weight $\mle{W}(r')$ & $63\nflock$ & $64\nflock$\\
\end{tabular}
\end{center}

Everything above is in $\E$. The $64$ products $x^i\,\Phi(s_i)$ in the target are left out: $x^i$ lies in $\K$, so they are cheaper, and stepping $x^i$ to $x^{i+1}$ is a shift.

\subsection{Credits}\label{rs:credits}

The construction above rests on generalizing ring switching~\cite{DP26} to extension degrees that are not powers of two, an idea of Lev Soukhanov, \textit{[[alloc]init]}. The map $\Phi$ follows~\cite{Flock26}.
